\begin{quote}
\noindent\textbf{Abstrak}---Model \textit{Mean-Variance} yang digagas oleh Markowitz merupakan model matematika untuk strategi optimasi portofolio investasi dengan metode pencarian alokasi asset yang memaksimalkan ekspektasi imbal hasil sekaligus meminimalkan varians sebagai risiko. Model tersebut membutuhkan algoritma optimasi yang mumpuni untuk menghasilkan portofolio optimal dari data historis harga pasar. Namun, model optimasi konvensional seringkali mengalami kendala dalam menangkap dinamika pasar modal yang kompleks, terutama saat menangani korelasi antar aset dalam jumlah besar. Untuk mengatasi hal tersebut, tugas akhir ini bertujuan untuk mengoptimalkan alokasi portofolio investasi berdasarkan model \textit{Mean-Variance} yang ditingkatkan lewat kerangka kerja optimasi portofolio hibrida dengan mengintegrasikan prinsip teori permainan potensial eksak ke dalam algoritma \textit{Variational Quantum Eigensolver} (VQE). Kombinasi ini dilakukan dengan pemetaan permasalahan alokasi portofolio ke dalam bentuk Hamiltonian Ising lewat penyetaraan yang diturunkan dari transformasi QUBO. Untuk mempercepat konvergensi proses optimasi pada komputer kuantum, algoritma pencarian ekuilibrium Nash digunakan sebagai inisialisasi keadaan awal alih-alih menggunakan inisialisasi acak. Hasil simulasi menggunakan data saham indeks IHSG periode 2021--2023 menunjukkan bahwa strategi hibrida GT-VQE secara keseluruhan mengungguli metode klasik dan kuantum murni dengan capaian rasio Sharpe sebesar 1,0132 pada sistem empat aset. Integrasi teori permainan terbukti meningkatkan kestabilan pertumbuhan modal di tengah volatilitas pasar serta mengefisiensikan kedalaman sirkuit kuantum. Temuan ini mengukuhkan potensi pendekatan kuantum-klasik hibrida sebagai solusi praktis yang tangguh untuk optimasi portofolio pada era NISQ. \\

\noindent\textbf{Kata Kunci}---\textit{Ekonofisika, Game Theory, Hamiltonian Ising, Optimasi Portofolio, VQE}
\end{quote}
